\section{The separation theorem}
\label{sec:synthesis}

The irrationality input in this section is now unconditional:
Theorem~\ref{thm:every-cubic} applies to every smooth cubic threefold.  The
universal \(CH_0\)-triviality inputs are independent classical or
cycle-theoretic statements, and pass to the product by the projective-bundle
formula for Chow groups.

\begin{proof}[Proof of Corollary~\ref{cor:voisin-separation}]
Voisin produces a nonempty countable union of subvarieties of codimension at
most three in the moduli space of smooth cubic threefolds along which the
primitive minimal class \(\Theta^4/4!\) is algebraic, hence along which
\(CH_0\) is universally trivial \cite[Theorem~4.5 and Lemma~4.6]{Voisin}.
The projective-bundle formula supplies universal \(CH_0\)-triviality of the
product, and Theorem~\ref{thm:every-cubic} supplies its irrationality.
\proves{cor:voisin-separation}%
\uses{thm:every-cubic}%
\imports{Voisin:minimal-class-locus}%
\end{proof}

\begin{proof}[Proof of Corollary~\ref{cor:fermat-separation}]
The Fermat equation is a sum of five forms in separated single variables, so
\cite{CT} gives universal \(CH_0\)-triviality; the projective-bundle formula
passes it to the product.  Irrationality follows from
Theorem~\ref{thm:every-cubic}.
\proves{cor:fermat-separation}%
\uses{thm:every-cubic}%
\imports{CT:separated-variable-criterion}%
\end{proof}

\begin{proof}[Proof of Corollary~\ref{cor:coprime-separation}]
The cubics of \cite[Theorem~3.3]{YYZ} carry unirational parametrizations of
degrees two and three, and \cite[Corollary~3.5]{YYZ} keeps both degrees on
\(X\times\PP^1\).  Parametrizations of coprime degrees give a decomposition
of the diagonal and hence universal \(CH_0\)-triviality
\cite[Section~1]{YYZ}.  Theorem~\ref{thm:every-cubic} gives irrationality of
\(X\times\PP^1\).
\proves{cor:coprime-separation}%
\uses{thm:every-cubic}%
\imports{YYZ:coprime-degree-family, YYZ:coprime-degrees-persist}%
\end{proof}

\begin{proof}[Proof of Theorem~\ref{thm:separation-family}]
Corollary~\ref{cor:universal-ch0} gives universal \(CH_0\)-triviality of
\(X_b\), and the projective-bundle formula gives the same property for
\(X_b\times\PP^1\).  Theorem~\ref{thm:every-cubic} proves that this
fourfold is irrational.  Non-isotriviality follows from the nonconstant
intermediate-Jacobian period map on the \(A_5\)-component: Hartlieb proves
that the closure of its image is a one-dimensional special subvariety
\cite[Proposition~5.7]{Hartlieb}.  Finally,
Proposition~\ref{prop:A5-nonseparated} shows that every point of this moduli
curve except the Fermat point lies outside the separated-variable locus.
\proves{thm:separation-family}%
\uses{cor:universal-ch0, thm:every-cubic, prop:A5-nonseparated}%
\imports{Hartlieb:period-image-dimension}%
\end{proof}

The two mechanisms are independent without any standing hypothesis.  On the
cycle side, algebraicity of \(\Theta^4/4!\) removes the
universal-\(CH_0\) obstruction.  On the quantum side, the rank-two residue
marker has value \(1\) on a cubic threefold and vanishes on every possible
center in a factorization of fourfolds; the categorical ledger therefore
obstructs rationality after one product stabilization.  The special
\(A_5\)-geometry enters only in the cycle argument, whereas the QDM marker
applies to every smooth cubic threefold.  The cycle construction also shows,
by Proposition~\ref{prop:A5-nonseparated}, that away from the Fermat point its
universal \(CH_0\)-triviality is not explained by the separated-variable
criterion.

The framed invariant \(\nu_6\) is a finer, partially conditional fold of the
same categorical ledger.  Its values \(\nu_6(X)=2\) and
\(\nu_6(X\times\PP^1)=4\) are unconditional for every smooth cubic
threefold by Proposition~\ref{prop:cubic-packet} and
Corollary~\ref{cor:cubic-product-nu}.  Hypothesis~5.7R is used only for the
residual reconstruction-tail step in the operation formulas.  The direct
specialization arguments of Propositions~\ref{prop:direct-specialized-lowdim}
and~\ref{prop:hirzebruch-specialized-vanishing}
remove Hypothesis~5.7T for nef-canonical targets, \(\PP^1\), \(\PP^2\),
every rational geometrically ruled surface (for \(F_1\), on the
specializations that arise as blowup centers), and, assuming 5.7R, ruled
surfaces over positive-genus curves; 5.7T remains only for surface centers that
are neither minimal nor geometrically ruled.
Conditional on 5.7R and on 5.7T for those centers, the common marker theorem
makes \(\nu_6\) birationally invariant through dimension four.  Neither
hypothesis is a condition on any separation statement in this section.

For comparison, the two threefold-level quantities are logically
independent.  If \(U(Y)\) denotes the indicator of universal
\(CH_0\)-triviality, Corollary~\ref{cor:p3-nu6} gives
\(\nu_6(\PP^3)=0\) unconditionally, while
Proposition~\ref{prop:cubic-packet} gives \(\nu_6=2\) for every cubic.
Consequently
\[
 \bigl(U(Y),\nu_6(Y)\bigr)=
 \begin{cases}
  (1,0),&Y=\PP^3,\\
  (1,2),&Y\text{ is a smooth member of the }A_5\text{-family},\\
  (0,2),&Y\text{ is a very general cubic threefold}.
 \end{cases}
\]
Voisin's criterion \cite[Corollary~4.4]{Voisin} and the evenness theorem
\cite[Theorem~1.3 and Corollary~1.4]{EdGFS} give the last line.  Thus neither
universal \(CH_0\)-triviality nor \(\nu_6\) determines the other.

All stabilization statements in this paper concern \(X\times\PP^1\).  No
claim about further stabilization is implicit in the theorem above.
